\documentclass[10pt,letterpaper]{article}

% Idioma, codificacion y tipografia
\usepackage[utf8]{inputenc}
\usepackage[T1]{fontenc}
\usepackage[spanish,es-nodecimaldot]{babel}
\usepackage{lmodern}
\usepackage{microtype}
\linespread{1.10}

% Pagina y composicion
\usepackage{geometry}
\geometry{
  letterpaper,
  top=1.8cm,
  bottom=1.8cm,
  left=1.8cm,
  right=1.8cm,
  headheight=19pt,
  headsep=0.45cm,
  footskip=0.75cm
}
\usepackage{multicol}
\usepackage{etoolbox}
\setlength{\columnsep}{0.58cm}
\setlength{\columnseprule}{0pt}
\setlength{\multicolsep}{0.55\baselineskip}
\setlength{\premulticols}{6\baselineskip}
\setlength{\postmulticols}{0.35\baselineskip}
\raggedcolumns

\clubpenalty=10000
\widowpenalty=10000
\displaywidowpenalty=10000
\setlength{\emergencystretch}{1.5em}
\setlength{\parindent}{0pt}
\setlength{\parskip}{0.38em}

% Color, tablas y listas
\usepackage[table]{xcolor}
\usepackage{booktabs}
\usepackage{longtable}
\usepackage{tabularx}
\usepackage{array}
\usepackage{enumitem}
\setlist[itemize]{leftmargin=1.35em,topsep=0.2em,itemsep=0.1em,parsep=0pt}
\setlist[enumerate]{leftmargin=1.55em,topsep=0.2em,itemsep=0.12em,parsep=0pt}

% Navegacion y estilo editorial
\usepackage{xurl}
\usepackage{hyperref}
\usepackage{fancyhdr}
\usepackage{titlesec}
\usepackage{listings}
\usepackage[most]{tcolorbox}

\definecolor{azulsql}{HTML}{123B5D}
\definecolor{azulmedio}{HTML}{176B87}
\definecolor{turquesa}{HTML}{2A9D8F}
\definecolor{naranjasql}{HTML}{E76F51}
\definecolor{tintasql}{HTML}{243746}
\definecolor{fondosuave}{HTML}{F3F7F9}
\definecolor{fondocodigo}{HTML}{F6F8FA}
\definecolor{bordecodigo}{HTML}{CAD6DE}
\definecolor{comentario}{HTML}{527A28}
\definecolor{cadena}{HTML}{9B2C53}
\definecolor{alertafondo}{HTML}{FFF7ED}
\definecolor{alertaborde}{HTML}{F59E0B}
\definecolor{errorfondo}{HTML}{FFF1F0}
\definecolor{errorborde}{HTML}{C2413B}

\hypersetup{
  colorlinks=true,
  linkcolor=azulsql,
  urlcolor=azulsql,
  pdftitle={Glosario esencial de SQL -- Unidad 1},
  pdfauthor={INF-324 Bases de Datos}
}

\pagestyle{fancy}
\fancyhf{}
\lhead{\fontsize{8.5}{10}\selectfont
  \textcolor{azulsql}{\textbf{INF-324 · Bases de Datos}}}
\rhead{\fontsize{8.5}{10}\selectfont
  \textcolor{tintasql}{Glosario SQL · Unidad 1}}
\cfoot{\fontsize{8.5}{10}\selectfont
  \textcolor{tintasql}{Página \thepage}}
\renewcommand{\headrulewidth}{0.35pt}
\renewcommand{\headrule}{\hbox to\headwidth{\color{bordecodigo}\leaders\hrule height \headrulewidth\hfill}}

\titleformat{\section}
  {\large\bfseries\color{azulsql}}
  {\colorbox{azulsql}{\textcolor{white}{\strut\thesection}}}{0.55em}{}
\titlespacing*{\section}{0pt}{1.15em}{0.45em}
\titleformat{\subsection}
  {\normalsize\bfseries\color{azulmedio}}
  {\thesubsection}{0.5em}{}
\titlespacing*{\subsection}{0pt}{0.85em}{0.25em}
\titleformat{\subsubsection}
  {\normalsize\bfseries\color{tintasql}}
  {\thesubsubsection}{0.45em}{}
\titlespacing*{\subsubsection}{0pt}{0.65em}{0.2em}

\lstdefinestyle{sql}{
  language=SQL,
  backgroundcolor=\color{fondocodigo},
  frame=single,
  rulecolor=\color{bordecodigo},
  framerule=0.55pt,
  framesep=0.35em,
  basicstyle=\ttfamily\fontsize{10}{10.8}\selectfont\color{tintasql},
  keywordstyle=\color{azulsql}\bfseries,
  commentstyle=\color{comentario}\itshape,
  stringstyle=\color{cadena},
  showstringspaces=false,
  breaklines=true,
  columns=fullflexible,
  keepspaces=true,
  tabsize=2,
  aboveskip=0.3em,
  belowskip=0.35em,
  xleftmargin=0.8em,
  xrightmargin=0.8em,
  framexleftmargin=0.65em,
  framexrightmargin=0.65em,
  literate=
    {á}{{\'a}}1 {é}{{\'e}}1 {í}{{\'i}}1 {ó}{{\'o}}1 {ú}{{\'u}}1
    {Á}{{\'A}}1 {É}{{\'E}}1 {Í}{{\'I}}1 {Ó}{{\'O}}1 {Ú}{{\'U}}1
    {ñ}{{\~n}}1 {Ñ}{{\~N}}1
}
\lstset{style=sql}

\newcommand{\clause}[1]{\textcolor{azulsql}{\textbf{#1}}}
\newcommand{\syntax}[1]{%
  \par\noindent\begingroup\raggedright
  \textbf{Sintaxis: }\texttt{\detokenize{#1}}\par
  \endgroup}

% Las ideas clave y objetivos se distinguen tipograficamente, sin caja.
\newenvironment{ideaclave}[1][Idea clave]
  {\par\smallskip\begingroup\raggedright
   \noindent\textcolor{turquesa}{\textbf{#1.}}\enspace}
  {\par\endgroup\smallskip}

\newenvironment{objetivoleccion}
  {\par\smallskip\noindent\color{tintasql}\small}
  {\par\medskip}

\newtcolorbox{precaucion}[1][Atención]{
  enhanced,breakable,colback=alertafondo,colframe=alertaborde,boxrule=0.55pt,
  arc=1.2mm,left=0.8em,right=0.8em,top=0.55em,bottom=0.55em,
  before skip=0.65em,after skip=0.65em,
  fonttitle=\bfseries,coltitle=white,colbacktitle=alertaborde,title=#1}

\newtcolorbox{errorcritico}[1][Error crítico]{
  enhanced,breakable,colback=errorfondo,colframe=errorborde,boxrule=0.55pt,
  arc=1.2mm,left=0.8em,right=0.8em,top=0.55em,bottom=0.55em,
  before skip=0.65em,after skip=0.65em,
  fonttitle=\bfseries,coltitle=white,colbacktitle=errorborde,title=#1}

% Caja reservada exclusivamente para ejemplos ejecutables seleccionados.
\newtcblisting{sqlexample}[1][Ejemplo ejecutable]{
  enhanced,listing only,
  colback=fondocodigo,colframe=bordecodigo,boxrule=0.55pt,arc=1.2mm,
  width=\textwidth,center,
  left=0.8em,right=0.8em,top=0.4em,bottom=0.4em,
  before skip=0.45em,after skip=0.5em,
  title={#1},fonttitle=\bfseries\small,coltitle=white,colbacktitle=azulsql,
  listing options={
    style=sql,
    frame=none,
    xleftmargin=0pt,
    xrightmargin=0pt,
    framexleftmargin=0pt,
    framexrightmargin=0pt,
    aboveskip=0pt,
    belowskip=0pt}}

\newcommand{\important}[1]{\begin{precaucion}#1\end{precaucion}}

\renewcommand{\arraystretch}{1.14}
\setlength{\tabcolsep}{5pt}
\newcommand{\tablerefsize}{\fontsize{9}{10.4}\selectfont}
\AtBeginEnvironment{tabular}{\tablerefsize}
\AtBeginEnvironment{tabularx}{\tablerefsize}
\AtBeginEnvironment{longtable}{\tablerefsize}

% Estado de la composicion teorica. Los bloques extensos suspenden las dos
% columnas y las restauran despues, por lo que nunca quedan partidos entre
% columnas.
\newif\ifintheorycolumns
\newcommand{\BeginTheory}{%
  \ifintheorycolumns\else
    \begin{multicols}{2}%
    \global\intheorycolumnstrue
  \fi}
\newcommand{\EndTheory}{%
  \ifintheorycolumns
    \global\intheorycolumnsfalse
    \end{multicols}%
  \fi}

\newif\ifresumeafterlisting
\BeforeBeginEnvironment{lstlisting}{%
  \ifintheorycolumns
    \global\resumeafterlistingtrue\EndTheory
  \else
    \global\resumeafterlistingfalse
  \fi
  \par\noindent\begin{minipage}{\textwidth}}
\AfterEndEnvironment{lstlisting}{%
  \end{minipage}\par
  \ifresumeafterlisting
    \global\resumeafterlistingfalse\BeginTheory
  \fi}

\newif\ifresumeafterexample
\BeforeBeginEnvironment{sqlexample}{%
  \ifintheorycolumns
    \global\resumeafterexampletrue\EndTheory
  \else
    \global\resumeafterexamplefalse
  \fi}
\AfterEndEnvironment{sqlexample}{%
  \ifresumeafterexample
    \global\resumeafterexamplefalse\BeginTheory
  \fi}

\newif\ifresumeafterlongtable
\BeforeBeginEnvironment{longtable}{%
  \ifintheorycolumns
    \global\resumeafterlongtabletrue\EndTheory
  \else
    \global\resumeafterlongtablefalse
  \fi}
\AfterEndEnvironment{longtable}{%
  \ifresumeafterlongtable
    \global\resumeafterlongtablefalse\BeginTheory
  \fi}

\newif\ifresumeaftertabularx
\BeforeBeginEnvironment{tabularx}{%
  \ifintheorycolumns
    \global\resumeaftertabularxtrue\EndTheory
  \else
    \global\resumeaftertabularxfalse
  \fi}
\AfterEndEnvironment{tabularx}{%
  \ifresumeaftertabularx
    \global\resumeaftertabularxfalse\BeginTheory
  \fi}

\newif\ifresumeaftercenter
\BeforeBeginEnvironment{center}{%
  \ifintheorycolumns
    \global\resumeaftercentertrue\EndTheory
  \else
    \global\resumeaftercenterfalse
  \fi}
\AfterEndEnvironment{center}{%
  \ifresumeaftercenter
    \global\resumeaftercenterfalse\BeginTheory
  \fi}

\newcommand{\indiceendoscolumnas}{%
  \begin{multicols}{2}
    \tableofcontents
  \end{multicols}}

\newenvironment{glossarylist}
  {\begin{description}[
    style=unboxed,
    leftmargin=0pt,
    labelindent=0pt,
    labelsep=0.45em,
    itemsep=0.18em,
    parsep=0pt,
    topsep=0.2em,
    font=\normalfont\ttfamily\bfseries]}
  {\end{description}}

\begin{document}
\hypersetup{pageanchor=false}

\begin{titlepage}
\thispagestyle{empty}
\color{tintasql}
\noindent\colorbox{azulsql}{%
  \parbox[c][1.05cm][c]{\dimexpr\textwidth-2\fboxsep\relax}{%
    \centering\color{white}\bfseries\large INF-324 · BASES DE DATOS}}

\vspace{2.3cm}
{\color{turquesa}\rule{2.2cm}{4pt}}\\[0.75cm]
{\fontsize{31}{36}\selectfont\bfseries Glosario esencial\\de SQL\par}
\vspace{0.45cm}
{\Large\color{azulmedio}Unidad 1 · Lecciones 1 a 6\par}
\vspace{0.35cm}
{\large Oracle Database · esquema OEHR\par}

\vfill
\begin{tcolorbox}[colback=fondosuave,colframe=bordecodigo,boxrule=0.5pt,
  arc=2mm,left=1.1em,right=1.1em,top=0.9em,bottom=0.9em]
\textbf{Propósito de esta guía}\par
Comprender, escribir y revisar las consultas estudiadas en la Unidad 1. Cada
tema conecta el propósito del comando con su sintaxis, un ejemplo ejecutable y
las precauciones que evitan los errores más frecuentes.
\end{tcolorbox}

\vspace{1.2cm}
\begin{tabularx}{\textwidth}{@{}X r@{}}
\textbf{Material de apoyo académico} & \textbf{Edición 2026} \\
Lecciones 1 a 6 & Consulta y práctica \\
\end{tabularx}
\vspace{0.35cm}
{\color{naranjasql}\rule{\textwidth}{2pt}}
\end{titlepage}
\hypersetup{pageanchor=true}
\pagenumbering{roman}
\BeginTheory

\section*{Cómo utilizar esta guía}

Esta guía reúne las cláusulas, operadores y funciones trabajados en las seis
lecciones de la Unidad 1. Los ejemplos están escritos para Oracle Database y
el conjunto de datos OEHR. Puede utilizarse de dos maneras:

\begin{itemize}
  \item \textbf{Lectura progresiva:} seguir las lecciones en orden para construir
        consultas cada vez más completas.
  \item \textbf{Consulta rápida:} buscar un término en el índice o en el
        glosario alfabético final.
\end{itemize}

\begin{ideaclave}[Cómo leer cada tema]
Primero identifique \textbf{qué problema resuelve} el comando; luego observe su
\textbf{sintaxis}; finalmente, ejecute el \textbf{ejemplo} y cambie una condición
para comprobar cómo varía el resultado. Las cajas naranjas señalan errores que
conviene anticipar.
\end{ideaclave}

\begin{precaucion}[Convención de tablas]
Los materiales del repositorio usan nombres como
\texttt{OEHR\_EMPLOYEES}. Si se trabaja con el esquema estándar HR, la tabla
equivalente se escribe \texttt{HR.EMPLOYEES}. Elija una convención y manténgala
en toda la consulta; no mezcle ambos nombres.
\end{precaucion}

\rowcolors{2}{fondosuave}{white}
\begin{center}
\begin{tabular}{ll}
\rowcolor{azulsql}
\textcolor{white}{\textbf{Tablas del repositorio}} &
\textcolor{white}{\textbf{Equivalente en el esquema HR}} \\
\texttt{OEHR\_EMPLOYEES}   & \texttt{HR.EMPLOYEES} \\
\texttt{OEHR\_DEPARTMENTS} & \texttt{HR.DEPARTMENTS} \\
\texttt{OEHR\_JOBS}        & \texttt{HR.JOBS} \\
\texttt{OEHR\_LOCATIONS}   & \texttt{HR.LOCATIONS} \\
\texttt{OEHR\_COUNTRIES}   & \texttt{HR.COUNTRIES} \\
\texttt{OEHR\_REGIONS}     & \texttt{HR.REGIONS} \\
\end{tabular}
\end{center}
\rowcolors{2}{}{}

\EndTheory
\clearpage
\indiceendoscolumnas
\clearpage
\pagenumbering{arabic}
\BeginTheory

\section{Mapa rápido de la unidad}

Antes de estudiar comandos aislados, conviene reconocer la progresión de la
unidad: primero se recuperan filas, luego se transforman y agrupan y, por
último, se relacionan varias tablas.

\rowcolors{2}{fondosuave}{white}
\begin{tabularx}{\textwidth}{>{\bfseries}p{1.5cm} X}
\rowcolor{azulsql}
\textcolor{white}{Lección} & \textcolor{white}{Contenidos principales} \\
1 & \lstinline|SELECT|, \lstinline|FROM|, \lstinline|DESCRIBE|, expresiones
aritméticas, alias, concatenación y \lstinline|DISTINCT|. \\
2 & \lstinline|WHERE|, operadores de comparación y lógicos,
\lstinline|BETWEEN|, \lstinline|IN|, \lstinline|LIKE|,
\lstinline|IS NULL| y \lstinline|ORDER BY|. \\
3 & Funciones de texto, funciones numéricas, tabla \lstinline|DUAL|,
\lstinline|SYSDATE|, aritmética y funciones de fecha. \\
4 & Conversión de tipos, formatos de fecha y número, funciones para nulos y
expresiones condicionales. \\
5 & Funciones de grupo, \lstinline|GROUP BY| y \lstinline|HAVING|. \\
6 & Alias de tabla, productos cartesianos y combinaciones con
\lstinline|JOIN|, \lstinline|ON| y \lstinline|USING|. \\
\end{tabularx}
\rowcolors{2}{}{}

\section{Anatomía de una consulta}

Cuando las cláusulas se utilizan juntas, su orden de escritura es estricto:

\begin{lstlisting}
SELECT   columnas o expresiones
FROM     tabla principal
JOIN     otra tabla
  ON     condicion de union
WHERE    condiciones sobre filas
GROUP BY columnas de agrupacion
HAVING   condiciones sobre grupos
ORDER BY columnas de ordenamiento;
\end{lstlisting}

\begin{ideaclave}[Regla de lectura]
Lea una consulta en dos direcciones: \textbf{de arriba abajo} para reconocer lo
que está escrito y desde \textbf{FROM/JOIN hacia SELECT} para comprender cómo se
construye el resultado. Esta segunda lectura explica qué datos existen en cada
etapa.
\end{ideaclave}

No todas son obligatorias. En una consulta común, \lstinline|SELECT| indica qué
se mostrará y \lstinline|FROM| indica de dónde se obtendrán los datos. Las otras
cláusulas se agregan solo cuando son necesarias. Puede haber varios
\lstinline|JOIN|, cada uno con su propia condición.

\begin{sqlexample}[Consulta completa: selección, unión, filtro y orden]
SELECT e.employee_id,
       e.first_name,
       d.department_name
FROM oehr_employees e
JOIN oehr_departments d
  ON e.department_id = d.department_id
WHERE d.department_name = 'IT'
ORDER BY e.last_name;
\end{sqlexample}

\subsection{Orden de escritura y orden lógico}

El orden en que Oracle procesa conceptualmente una consulta no es el mismo en
que se escribe. Esta diferencia ayuda a comprender por qué, por ejemplo,
\lstinline|WHERE| no puede usar una función de grupo.

\begin{center}
\begin{tabular}{cll}
\toprule
\rowcolor{azulsql}
\textcolor{white}{\textbf{Paso}} & \textcolor{white}{\textbf{Escritura}} &
\textcolor{white}{\textbf{Procesamiento lógico simplificado}} \\
\midrule
1 & \lstinline|SELECT| & \lstinline|FROM| y \lstinline|JOIN| \\
2 & \lstinline|FROM| & \lstinline|WHERE| \\
3 & \lstinline|JOIN ... ON| & \lstinline|GROUP BY| \\
4 & \lstinline|WHERE| & \lstinline|HAVING| \\
5 & \lstinline|GROUP BY| & \lstinline|SELECT| \\
6 & \lstinline|HAVING| & \lstinline|ORDER BY| \\
7 & \lstinline|ORDER BY| & --- \\
\bottomrule
\end{tabular}
\end{center}

\section{Lección 1: recuperación básica de datos}

\begin{objetivoleccion}
\textbf{Objetivo de aprendizaje.} Recuperar columnas de una tabla, calcular
expresiones y presentar encabezados y textos comprensibles sin modificar los
datos almacenados.
\end{objetivoleccion}

\subsection{SELECT y FROM}

\clause{SELECT} elige columnas o calcula expresiones. \clause{FROM} señala la
tabla que contiene los datos.

\syntax{SELECT columna_1, columna_2 FROM nombre_tabla;}

\begin{lstlisting}
SELECT employee_id, first_name, salary
FROM oehr_employees;
\end{lstlisting}

El asterisco obtiene todas las columnas. Es útil para explorar, pero en reportes
conviene escribir únicamente las columnas necesarias.

\begin{lstlisting}
SELECT *
FROM oehr_departments;
\end{lstlisting}

\subsection{DESCRIBE o DESC}

Muestra la estructura de una tabla: columnas, tipos de datos y posibilidad de
valores nulos. Es un comando del cliente de Oracle, no una consulta
\lstinline|SELECT|.

\syntax{DESCRIBE nombre_tabla;}

\begin{lstlisting}
DESCRIBE oehr_locations;
-- Forma abreviada:
DESC oehr_locations;
\end{lstlisting}

\subsection{Expresiones aritméticas}

Se pueden usar \lstinline|+|, \lstinline|-|, \lstinline|*| y
\lstinline|/| dentro de \lstinline|SELECT|. La multiplicación y la división se
evalúan antes que la suma y la resta; los paréntesis permiten controlar el
orden.

\syntax{SELECT expresion_aritmetica AS alias FROM nombre_tabla;}

\begin{sqlexample}[Ejemplo ejecutable: expresiones aritméticas]
SELECT employee_id,
       salary,
       salary * 12 AS salario_anual,
       (salary + 57) * 12 AS anual_con_bono_mensual,
       (salary * 12) + 57 AS anual_con_bono_anual
FROM oehr_employees;
\end{sqlexample}

Si una expresión opera con \lstinline|NULL|, su resultado también será
\lstinline|NULL|. Más adelante se puede usar \lstinline|NVL| para sustituirlo.

\subsection{Alias de columna}

Un alias cambia el encabezado mostrado, pero no modifica la tabla. Las comillas
dobles permiten conservar espacios, acentos y mayúsculas exactas.

\syntax{expresion AS "Encabezado visible"}

\begin{lstlisting}
SELECT employee_id AS "COD_EMPLEADO",
       first_name AS "Nombre",
       salary * 12 AS "Sueldo anual"
FROM oehr_employees;
\end{lstlisting}

\subsection{Concatenación y texto literal}

El operador \lstinline@||@ une valores de texto. Los textos literales se
escriben entre comillas simples.

\syntax{expresion_1 || ' texto ' || expresion_2}

\begin{lstlisting}
SELECT employee_id,
       first_name || ' ' || last_name AS nombre_completo,
       'El empleado ' || first_name || ' tiene el cargo ' || job_id AS reporte
FROM oehr_employees;
\end{lstlisting}

\subsection{DISTINCT}

Elimina filas duplicadas del resultado según el conjunto de columnas elegido.

\syntax{SELECT DISTINCT columna_1, columna_2 FROM nombre_tabla;}

\begin{lstlisting}
SELECT DISTINCT job_id
FROM oehr_employees;
\end{lstlisting}

\section{Lección 2: filtrar y ordenar}

\begin{objetivoleccion}
\textbf{Objetivo de aprendizaje.} Restringir filas mediante condiciones,
combinar criterios de búsqueda y controlar el orden del resultado final.
\end{objetivoleccion}

\subsection{WHERE y operadores de comparación}

\clause{WHERE} conserva únicamente las filas que cumplen una condición.

\syntax{SELECT columnas FROM tabla WHERE condicion;}

\begin{center}
\begin{tabular}{cl@{\hspace{2cm}}cl}
\toprule
\rowcolor{azulsql}
\textcolor{white}{\textbf{Operador}} & \textcolor{white}{\textbf{Significado}} &
\textcolor{white}{\textbf{Operador}} & \textcolor{white}{\textbf{Significado}} \\
\midrule
\lstinline|=|  & igual & \lstinline|<>| o \lstinline|!=| & distinto \\
\lstinline|>|  & mayor que & \lstinline|>=| & mayor o igual \\
\lstinline|<|  & menor que & \lstinline|<=| & menor o igual \\
\bottomrule
\end{tabular}
\end{center}

\begin{lstlisting}
SELECT employee_id, last_name, salary
FROM oehr_employees
WHERE salary >= 5000;
\end{lstlisting}

Las cadenas distinguen mayúsculas y minúsculas y se escriben entre comillas
simples. Para fechas es preferible un literal ANSI, pues no depende del formato
regional configurado en la sesión.

\begin{lstlisting}
SELECT employee_id, last_name, hire_date
FROM oehr_employees
WHERE hire_date >= DATE '2016-01-01';
\end{lstlisting}

\subsection{AND, OR y NOT}

\lstinline|AND| exige que ambas condiciones sean verdaderas;
\lstinline|OR| exige al menos una; \lstinline|NOT| niega una condición. Oracle
evalúa \lstinline|NOT|, después \lstinline|AND| y finalmente \lstinline|OR|.
Los paréntesis hacen explícita la intención.

\syntax{WHERE (condicion_1 OR condicion_2) AND condicion_3}

\begin{sqlexample}[Ejemplo ejecutable: condiciones combinadas]
SELECT employee_id, last_name, department_id, salary
FROM oehr_employees
WHERE (department_id = 50 OR department_id = 80)
  AND salary >= 5000;
\end{sqlexample}

\subsection{BETWEEN y NOT BETWEEN}

Busca dentro o fuera de un rango inclusivo: considera ambos límites.

\syntax{WHERE columna [NOT] BETWEEN limite_inferior AND limite_superior}

\begin{lstlisting}
SELECT employee_id, last_name, salary
FROM oehr_employees
WHERE salary BETWEEN 3000 AND 5000;
\end{lstlisting}

\begin{lstlisting}
SELECT employee_id, last_name, hire_date
FROM oehr_employees
WHERE hire_date NOT BETWEEN DATE '2015-01-01' AND DATE '2020-12-31';
\end{lstlisting}

\subsection{IN y NOT IN}

Compara una columna con una lista de valores. Es una forma clara de expresar
varias comparaciones unidas mediante \lstinline|OR|.

\syntax{WHERE columna [NOT] IN (valor_1, valor_2, valor_3)}

\begin{lstlisting}
SELECT employee_id, last_name, department_id
FROM oehr_employees
WHERE department_id IN (50, 80, 90);
\end{lstlisting}

\begin{lstlisting}
SELECT employee_id, last_name, manager_id
FROM oehr_employees
WHERE manager_id NOT IN (100, 114, 120, 121);
\end{lstlisting}

\important{Si la lista de \texttt{NOT IN} contiene un valor
\texttt{NULL}, el resultado puede no devolver filas debido a la lógica de
tres valores de SQL. En ese caso se debe revisar el origen de la lista.}

\subsection{LIKE y comodines}

\lstinline|LIKE| compara patrones de texto. El símbolo \lstinline|%| representa
cero o más caracteres y \lstinline|_| representa exactamente un carácter.

\syntax{WHERE columna [NOT] LIKE 'patron'}

\begin{center}
\begin{tabular}{ll}
\toprule
\rowcolor{azulsql}
\textcolor{white}{\textbf{Patrón}} & \textcolor{white}{\textbf{Significado}} \\
\midrule
\lstinline|'S%'| & comienza con S \\
\lstinline|'%a'| & termina con a \\
\lstinline|'%an%'| & contiene an \\
\lstinline|'_a%s'| & la segunda letra es a y termina con s \\
\bottomrule
\end{tabular}
\end{center}

\begin{lstlisting}
SELECT employee_id, last_name
FROM oehr_employees
WHERE last_name LIKE 'S%'
   OR last_name LIKE '%a';
\end{lstlisting}

Para buscar literalmente un guion bajo o un porcentaje se puede definir un
carácter de escape:

\begin{lstlisting}
SELECT employee_id, job_id
FROM oehr_employees
WHERE job_id LIKE '%SA\_%' ESCAPE '\';
\end{lstlisting}

\important{\texttt{= 'E\%'} busca exactamente los dos caracteres
\texttt{E\%}. Para buscar textos que comienzan con E se debe usar
\texttt{LIKE 'E\%'}. Aun así, el resultado será vacío si los datos unidos no
contienen un país que comience con esa letra.}

\subsection{IS NULL e IS NOT NULL}

\lstinline|NULL| representa un valor ausente o desconocido. No se compara con
\lstinline|=| ni con \lstinline|<>|.

\syntax{WHERE columna IS [NOT] NULL}

\begin{lstlisting}
SELECT employee_id, last_name, commission_pct
FROM oehr_employees
WHERE commission_pct IS NULL;
\end{lstlisting}

\begin{lstlisting}
SELECT employee_id, last_name, commission_pct
FROM oehr_employees
WHERE commission_pct IS NOT NULL;
\end{lstlisting}

\subsection{ORDER BY}

Ordena el resultado. \lstinline|ASC| es ascendente y es el valor por defecto;
\lstinline|DESC| es descendente. Se puede ordenar por una columna, un alias, una
expresión o la posición de una columna del \lstinline|SELECT|.

\syntax{ORDER BY columna_1 [ASC|DESC], columna_2 [ASC|DESC];}

\begin{lstlisting}
SELECT employee_id, last_name, department_id, hire_date
FROM oehr_employees
ORDER BY department_id ASC, hire_date DESC;
\end{lstlisting}

Aunque \lstinline|ORDER BY 2| es válido, usar el nombre o el alias suele ser más
claro y resiste mejor los cambios en la lista de columnas.

\section{Lección 3: funciones de una sola fila}

Estas funciones producen un resultado por cada fila de entrada.

\begin{objetivoleccion}
\textbf{Objetivo de aprendizaje.} Transformar textos, números y fechas fila a
fila, sin perder de vista el tipo de dato que recibe y devuelve cada función.
\end{objetivoleccion}

\begin{ideaclave}[Cómo interpretar una función]
En \texttt{FUNCION(argumento)}, Oracle evalúa el argumento de cada fila y
devuelve un valor. Las funciones pueden anidarse: la más interna se ejecuta
primero y su resultado alimenta a la siguiente.
\end{ideaclave}

\subsection{Funciones de texto}

\begin{longtable}{>{\ttfamily}p{3.1cm} p{5.7cm} p{6.1cm}}
\toprule
\rowcolor{azulsql}
\normalfont\textcolor{white}{\textbf{Función}} &
\textcolor{white}{\textbf{Propósito}} & \textcolor{white}{\textbf{Ejemplo}} \\
\midrule
\endhead
UPPER & Convierte a mayúsculas. & \lstinline|UPPER(last_name)| \\
LOWER & Convierte a minúsculas. & \lstinline|LOWER(email)| \\
INITCAP & Coloca en mayúscula la inicial de cada palabra. &
\lstinline|INITCAP(first_name)| \\
CONCAT & Une dos valores. Para más valores debe anidarse o usarse
\lstinline@||@. & \lstinline|CONCAT(first_name, last_name)| \\
LENGTH & Cuenta caracteres. & \lstinline|LENGTH(last_name)| \\
SUBSTR & Extrae una parte desde una posición y con una longitud opcional. &
\lstinline|SUBSTR(first_name, 1, 1)| \\
INSTR & Devuelve la posición de una cadena dentro de otra; cero si no existe. &
\lstinline|INSTR(last_name, 'a')| \\
TRIM & Elimina caracteres en los extremos; por defecto, espacios. &
\lstinline|TRIM('  texto  ')| \\
LPAD & Completa a la izquierda hasta una longitud. &
\lstinline|LPAD(employee_id, 6, '0')| \\
RPAD & Completa a la derecha hasta una longitud. &
\lstinline|RPAD(first_name, 20, '.')| \\
REPLACE & Sustituye apariciones de un texto. &
\lstinline|REPLACE(phone_number, '.', '-')| \\
\bottomrule
\end{longtable}

Ejemplo combinado:

\begin{sqlexample}[Ejemplo ejecutable: funciones de texto]
SELECT employee_id,
       UPPER(first_name || ' ' || last_name) AS nombre,
       LENGTH(last_name) AS letras,
       INSTR(LOWER(last_name), 'a') AS posicion_a
FROM oehr_employees
WHERE SUBSTR(last_name, -1, 1) = 'n';
\end{sqlexample}

\subsection{Funciones numéricas y DUAL}

\lstinline|ROUND(numero, n)| redondea; \lstinline|TRUNC(numero, n)| corta sin
redondear; \lstinline|MOD(a, b)| devuelve el resto de dividir \emph{a} por
\emph{b}. \lstinline|DUAL| es una tabla especial de Oracle útil para evaluar
expresiones sin consultar una tabla de negocio.

\syntax{ROUND(numero, decimales) | TRUNC(numero, decimales) | MOD(a, b)}

\begin{lstlisting}
SELECT ROUND(45.5674, 2) AS redondeado,
       TRUNC(45.5674, 2) AS truncado,
       MOD(54, 3) AS resto
FROM dual;
\end{lstlisting}

El separador decimal escrito en SQL suele ser el punto. La forma correcta de
\lstinline|MOD| recibe dos argumentos: \lstinline|MOD(54, 3)|.

\subsection{SYSDATE y aritmética de fechas}

\lstinline|SYSDATE| devuelve la fecha y hora actuales del servidor Oracle. No
lleva paréntesis.

\syntax{SELECT SYSDATE FROM DUAL;}

\begin{lstlisting}
SELECT SYSDATE
FROM dual;
\end{lstlisting}

En Oracle, sumar o restar un número a una fecha suma o resta días. La diferencia
entre dos fechas produce una cantidad de días.

\begin{lstlisting}
SELECT SYSDATE + 1 AS manana,
       SYSDATE - 7 AS hace_una_semana
FROM dual;
\end{lstlisting}

\begin{lstlisting}
SELECT employee_id,
       hire_date,
       TRUNC(SYSDATE - hire_date) AS dias_trabajados,
       TRUNC((SYSDATE - hire_date) / 7) AS semanas_trabajadas
FROM oehr_employees;
\end{lstlisting}

\important{Una columna \texttt{DATE} también almacena hora. Por eso
\texttt{fecha = SYSDATE} rara vez encuentra una fila: tendrían que coincidir
fecha y hora. Para comparar solo el día puede usarse
\texttt{TRUNC(fecha) = TRUNC(SYSDATE)}.}

\subsection{Funciones de fecha}

\begin{longtable}{>{\ttfamily}p{4.5cm} p{10.3cm}}
\toprule
\rowcolor{azulsql}
\normalfont\textcolor{white}{\textbf{Función}} &
\textcolor{white}{\textbf{Descripción}} \\
\midrule
\endhead
MONTHS\_BETWEEN(f1,f2) & Cantidad de meses entre dos fechas; devuelve un número. \\
ADD\_MONTHS(fecha,n) & Suma o resta meses a una fecha. \\
NEXT\_DAY(fecha,día) & Devuelve la siguiente ocurrencia del día indicado. El
idioma del nombre depende de \lstinline|NLS_DATE_LANGUAGE|. \\
LAST\_DAY(fecha) & Devuelve el último día del mes de la fecha. \\
ROUND(fecha,formato) & Redondea una fecha a la unidad indicada. \\
TRUNC(fecha,formato) & Trunca una fecha a la unidad indicada; sin formato quita
la hora. \\
\bottomrule
\end{longtable}

\begin{lstlisting}
SELECT employee_id,
       hire_date,
       MONTHS_BETWEEN(SYSDATE, hire_date) AS meses,
       ADD_MONTHS(hire_date, 6) AS revision,
       NEXT_DAY(hire_date, 'MONDAY') AS lunes_siguiente,
       LAST_DAY(hire_date) AS fin_de_mes
FROM oehr_employees;
\end{lstlisting}

\section{Lección 4: conversiones, nulos y condiciones}

\begin{objetivoleccion}
\textbf{Objetivo de aprendizaje.} Convertir datos de forma explícita, controlar
la presentación de fechas y números, tratar valores nulos y producir resultados
condicionales.
\end{objetivoleccion}

\subsection{TO\_CHAR para fechas}

Convierte una fecha a texto utilizando un modelo de formato.

\syntax{TO_CHAR(fecha, 'modelo_de_formato')}

\begin{center}
\begin{tabular}{ll@{\hspace{1.4cm}}ll}
\toprule
\rowcolor{azulsql}
\textcolor{white}{\textbf{Formato}} & \textcolor{white}{\textbf{Significado}} &
\textcolor{white}{\textbf{Formato}} & \textcolor{white}{\textbf{Significado}} \\
\midrule
\lstinline|DD| & día del mes & \lstinline|MM| & mes numérico \\
\lstinline|MONTH| & nombre del mes & \lstinline|YYYY| & año de cuatro dígitos \\
\lstinline|DAY| & nombre del día & \lstinline|HH24| & hora de 00 a 23 \\
\lstinline|MI| & minutos & \lstinline|SS| & segundos \\
\lstinline|MON| & mes abreviado & \lstinline|DY| & día abreviado \\
\lstinline|YEAR| & año escrito en letras & \lstinline|AM / PM| & indicador de meridiano \\
\bottomrule
\end{tabular}
\end{center}

\begin{lstlisting}
SELECT TO_CHAR(SYSDATE, 'DD/MM/YYYY HH24:MI:SS') AS fecha_hora,
       TO_CHAR(SYSDATE, 'DAY DD "DE" MONTH "DE" YYYY') AS fecha_larga
FROM dual;
\end{lstlisting}

El modificador \lstinline|FM| elimina espacios y ceros de relleno agregados por
algunos formatos. El texto escrito entre comillas dobles dentro del modelo se
reproduce literalmente.

\begin{lstlisting}
SELECT TO_CHAR(SYSDATE, 'FMDAY DD "DE" MONTH "DE" YYYY')
FROM dual;
\end{lstlisting}

\subsection{TO\_CHAR para números}

Convierte un número en texto con el formato de presentación solicitado.

\syntax{TO_CHAR(numero, 'modelo_de_formato')}

En los modelos numéricos, \lstinline|9| reserva una posición opcional,
\lstinline|0| obliga a mostrar un cero, \lstinline|,| y \lstinline|.| representan
separadores y \lstinline|$| imprime el símbolo indicado. La configuración NLS
de la sesión puede influir en algunos separadores y símbolos monetarios.

\begin{lstlisting}
SELECT employee_id,
       TO_CHAR(salary, '$999,999.99') AS sueldo_formateado
FROM oehr_employees;
\end{lstlisting}

El resultado de \lstinline|TO_CHAR| es texto. Sirve para presentar datos, no
para continuar cálculos numéricos.

\subsection{TO\_NUMBER y TO\_DATE}

\lstinline|TO_NUMBER| convierte texto a número y \lstinline|TO_DATE| convierte
texto a fecha usando el formato indicado.

\syntax{TO_NUMBER(texto, 'formato') | TO_DATE(texto, 'formato')}

\begin{lstlisting}
SELECT TO_NUMBER('2026', '9999') AS anio,
       TO_DATE('11/09/2026', 'DD/MM/YYYY') AS fecha
FROM dual;
\end{lstlisting}

Cuando una fecha es fija, el literal ANSI suele ser más claro:

\begin{lstlisting}
SELECT DATE '2026-09-11' AS fecha
FROM dual;
\end{lstlisting}

Se deben evitar conversiones implícitas como \lstinline|hire_date > '01/01/2016'|,
porque dependen de la configuración regional y pueden fallar o interpretarse de
forma distinta.

\subsection{NVL, NVL2, NULLIF y COALESCE}

Estas funciones permiten decidir qué valor devolver cuando un dato es nulo o
cuando dos expresiones resultan iguales.

\begin{longtable}{>{\ttfamily}p{4.3cm} p{10.4cm}}
\toprule
\rowcolor{azulsql}
\normalfont\textcolor{white}{\textbf{Función}} &
\textcolor{white}{\textbf{Resultado}} \\
\midrule
\endhead
NVL(valor,reemplazo) & Devuelve el reemplazo cuando el valor es
\lstinline|NULL|; en caso contrario devuelve el valor original. \\
NVL2(v,no\_nulo,nulo) & Elige entre dos resultados según el valor
sea o no nulo. \\
NULLIF(a,b) & Devuelve \lstinline|NULL| si \emph{a} y \emph{b} son iguales; si
no, devuelve \emph{a}. \\
COALESCE(a,b,c,...) & Devuelve el primer argumento que no sea
\lstinline|NULL|. \\
\bottomrule
\end{longtable}

\begin{lstlisting}
SELECT employee_id,
       NVL(commission_pct, 0) AS comision,
       NVL2(commission_pct, 'CON COMISION', 'SIN COMISION') AS estado
FROM oehr_employees;
\end{lstlisting}

\begin{lstlisting}
SELECT NULLIF(10, 10) AS resultado_nulo,
       COALESCE(NULL, NULL, 'Primer valor disponible') AS resultado
FROM dual;
\end{lstlisting}

Los posibles resultados de \lstinline|NVL| y las alternativas de
\lstinline|NVL2| deben ser de tipos compatibles.

\subsection{CASE y DECODE}

\lstinline|CASE| expresa condiciones de manera legible. \lstinline|DECODE| es
una función propia de Oracle para comparaciones por igualdad; cuando la lógica
es más compleja suele preferirse \lstinline|CASE|.

\syntax{CASE WHEN condicion THEN resultado ELSE alternativa END}

\begin{sqlexample}[Ejemplo ejecutable: expresión CASE]
SELECT employee_id,
       salary,
       CASE
         WHEN salary >= 10000 THEN 'ALTO'
         WHEN salary >= 5000  THEN 'MEDIO'
         ELSE 'BAJO'
       END AS nivel_sueldo
FROM oehr_employees;
\end{sqlexample}

\begin{lstlisting}
SELECT employee_id,
       DECODE(department_id,
              50, 'Shipping',
              80, 'Sales',
              'Otro') AS area
FROM oehr_employees;
\end{lstlisting}

\section{Lección 5: funciones de grupo}

Las funciones de grupo producen un resultado por conjunto de filas e ignoran
valores \lstinline|NULL|, excepto \lstinline|COUNT(*)|, que cuenta filas.

\begin{objetivoleccion}
\textbf{Objetivo de aprendizaje.} Resumir conjuntos de filas mediante totales,
promedios, extremos y conteos; después, formar y filtrar grupos correctamente.
\end{objetivoleccion}

\begin{center}
\begin{tabular}{ll}
\toprule
\rowcolor{azulsql}
\textcolor{white}{\textbf{Función}} & \textcolor{white}{\textbf{Propósito}} \\
\midrule
\lstinline|AVG(expresion)| & promedio \\
\lstinline|SUM(expresion)| & suma \\
\lstinline|MIN(expresion)| & mínimo \\
\lstinline|MAX(expresion)| & máximo \\
\lstinline|COUNT(expresion)| & cantidad de valores no nulos \\
\lstinline|COUNT(*)| & cantidad de filas \\
\lstinline|COUNT(DISTINCT expresion)| & cantidad de valores distintos no nulos \\
\bottomrule
\end{tabular}
\end{center}

\begin{lstlisting}
SELECT ROUND(AVG(salary), 0) AS promedio,
       SUM(salary) AS total,
       MIN(hire_date) AS primera_contratacion,
       MAX(hire_date) AS ultima_contratacion,
       COUNT(*) AS empleados
FROM oehr_employees;
\end{lstlisting}

\subsection{GROUP BY}

Forma grupos. Toda expresión del \lstinline|SELECT| que no esté dentro de una
función de grupo debe aparecer en \lstinline|GROUP BY|.

\syntax{SELECT columna, FUNCION_GRUPO(...) FROM tabla GROUP BY columna;}

\begin{lstlisting}
SELECT department_id,
       job_id,
       TO_CHAR(SUM(salary), '$999,999.99') AS total_sueldos
FROM oehr_employees
GROUP BY department_id, job_id
ORDER BY department_id;
\end{lstlisting}

\subsection{WHERE frente a HAVING}

\lstinline|WHERE| filtra filas antes de agrupar. \lstinline|HAVING| filtra los
grupos ya calculados y puede usar funciones de grupo.

\syntax{WHERE condicion_de_fila GROUP BY columnas HAVING condicion_de_grupo}

\begin{sqlexample}[Ejemplo ejecutable: agrupación y filtro de grupos]
SELECT department_id,
       MAX(salary) AS sueldo_maximo
FROM oehr_employees
WHERE department_id BETWEEN 50 AND 90
GROUP BY department_id
HAVING MAX(salary) > 10000
ORDER BY department_id DESC;
\end{sqlexample}

No es válido escribir \lstinline|WHERE MAX(salary) > 10000| porque el máximo
todavía no existe en la etapa donde se ejecuta \lstinline|WHERE|.

\section{Lección 6: consultas multitabla}

\begin{objetivoleccion}
\textbf{Objetivo de aprendizaje.} Relacionar tablas mediante sus claves,
calificar columnas con alias y elegir el tipo de combinación que conserva las
filas necesarias para el reporte.
\end{objetivoleccion}

\subsection{Idea central del JOIN}

Un \lstinline|JOIN| no modifica ni mezcla físicamente las tablas. Construye un
resultado temporal relacionando filas. La condición de \lstinline|ON| suele
comparar una clave foránea con la clave primaria correspondiente.

\syntax{FROM tabla_a a JOIN tabla_b b ON a.clave_foranea = b.clave_primaria}

\begin{lstlisting}
SELECT e.employee_id,
       e.first_name,
       d.department_name
FROM oehr_employees e
JOIN oehr_departments d
  ON e.department_id = d.department_id;
\end{lstlisting}

Los alias \lstinline|e| y \lstinline|d| hacen la consulta más corta y evitan
ambigüedad cuando dos tablas contienen columnas con el mismo nombre.

\subsection{Camino entre varias tablas}

Para obtener el país y la región de un empleado se sigue la relación del modelo:

\begin{center}
\texttt{EMPLOYEES} $\rightarrow$ \texttt{DEPARTMENTS} $\rightarrow$
\texttt{LOCATIONS} $\rightarrow$ \texttt{COUNTRIES} $\rightarrow$
\texttt{REGIONS}
\end{center}

\begin{sqlexample}[Ejemplo ejecutable: recorrido entre cinco tablas]
SELECT e.employee_id AS "ID empleado",
       e.first_name AS "Nombre",
       c.country_name AS "Pais",
       l.city AS "Ciudad",
       r.region_name AS "Region"
FROM oehr_employees e
JOIN oehr_departments d
  ON e.department_id = d.department_id
JOIN oehr_locations l
  ON d.location_id = l.location_id
JOIN oehr_countries c
  ON l.country_id = c.country_id
JOIN oehr_regions r
  ON c.region_id = r.region_id
ORDER BY e.last_name;
\end{sqlexample}

Para conocer qué países existen realmente en ese resultado antes de filtrar:

\begin{lstlisting}
SELECT DISTINCT c.country_name
FROM oehr_employees e
JOIN oehr_departments d
  ON e.department_id = d.department_id
JOIN oehr_locations l
  ON d.location_id = l.location_id
JOIN oehr_countries c
  ON l.country_id = c.country_id
ORDER BY c.country_name;
\end{lstlisting}

\subsection{ON y USING}

\lstinline|ON| permite escribir cualquier condición de unión y es la forma más
explícita. \lstinline|USING| se puede usar cuando la columna relacionada tiene
el mismo nombre en ambas tablas.

\syntax{tabla_a JOIN tabla_b USING (columna_comun)}

\begin{lstlisting}
SELECT l.city, d.department_name
FROM oehr_locations l
JOIN oehr_departments d USING (location_id);
\end{lstlisting}

La columna escrita dentro de \lstinline|USING| no debe llevar alias dentro de la
lista: se escribe \lstinline|USING (location_id)|, no
\lstinline|USING (l.location_id)|.

\subsection{Tipos de combinación}

\begin{longtable}{>{\ttfamily}p{3.5cm} p{11.2cm}}
\toprule
\rowcolor{azulsql}
\normalfont\textcolor{white}{\textbf{Combinación}} &
\textcolor{white}{\textbf{Comportamiento}} \\
\midrule
\endhead
INNER JOIN / JOIN & Devuelve solo filas con coincidencia en ambas tablas. \\
LEFT JOIN & Conserva todas las filas de la tabla izquierda; coloca
\lstinline|NULL| cuando no hay coincidencia a la derecha. \\
RIGHT JOIN & Conserva todas las filas de la tabla derecha. \\
FULL OUTER JOIN & Conserva las filas de ambos lados, coincidan o no. \\
CROSS JOIN & Genera todas las combinaciones posibles; es un producto
cartesiano. \\
NATURAL JOIN & Une automáticamente columnas del mismo nombre. Debe usarse con
cuidado porque una nueva columna homónima puede cambiar el resultado. \\
\bottomrule
\end{longtable}

Ejemplo para mostrar empleados aunque no tengan departamento:

\begin{lstlisting}
SELECT e.employee_id,
       e.last_name,
       d.department_name
FROM oehr_employees e
LEFT JOIN oehr_departments d
  ON e.department_id = d.department_id;
\end{lstlisting}

\subsection{Autocombinación y unión no igualitaria}

Una autocombinación usa dos alias de la misma tabla. Aquí se relaciona cada
empleado con su jefe:

\begin{lstlisting}
SELECT e.last_name AS empleado,
       j.last_name AS jefe
FROM oehr_employees e
LEFT JOIN oehr_employees j
  ON e.manager_id = j.employee_id;
\end{lstlisting}

Una unión no igualitaria usa una condición diferente de \lstinline|=|, por
ejemplo para ubicar un salario dentro de un rango:

\begin{lstlisting}
SELECT e.last_name, e.salary, g.grade_level
FROM oehr_employees e
JOIN job_grades g
  ON e.salary BETWEEN g.lowest_sal AND g.highest_sal;
\end{lstlisting}

La tabla \lstinline|JOB_GRADES| es ilustrativa y debe existir para ejecutar el
ejemplo.

\important{Si se enumeran tablas sin una condición correcta, se genera un
producto cartesiano: cada fila de una tabla se combina con cada fila de la otra.
La consulta puede ejecutarse, pero sus datos no representan la relación real.}

\section{DDL y DML observados brevemente}

En la parte 2 de la lección 3 aparecen dos instrucciones fuera de la consulta de
datos. Se registran aquí porque fueron observadas, aunque su estudio detallado
corresponde a otros bloques de la unidad.

\begin{objetivoleccion}
\textbf{Alcance.} En esta unidad se observaron \texttt{CREATE TABLE} e
\texttt{INSERT INTO}. Se incluyen para reconocer su propósito y su forma
básica; el estudio completo de DDL y DML se realiza en unidades posteriores.
\end{objetivoleccion}

\subsection{CREATE TABLE}

Crea una tabla y define sus columnas y tipos.

\syntax{CREATE TABLE nombre (columna tipo_de_dato);}

\begin{lstlisting}
CREATE TABLE ejemplo (
  rut VARCHAR2(12)
);
\end{lstlisting}

\subsection{INSERT INTO}

Agrega una fila. Los valores de texto deben ir entre comillas simples.

\syntax{INSERT INTO tabla (columnas) VALUES (valores);}

\begin{lstlisting}
INSERT INTO ejemplo (rut)
VALUES ('13.572.343-2');
\end{lstlisting}

\section{Errores frecuentes y cómo leerlos}

Un mensaje de error no es solo un obstáculo: indica en qué parte de la consulta
debe revisarse la sintaxis o la lógica. La siguiente tabla traduce los casos más
habituales a una acción concreta.

\begin{longtable}{>{\raggedright\arraybackslash}p{5cm}
                  >{\raggedright\arraybackslash}p{5cm}
                  >{\raggedright\arraybackslash}p{5cm}}
\toprule
\rowcolor{azulsql}
\textcolor{white}{\textbf{Situación}} &
\textcolor{white}{\textbf{Causa común}} &
\textcolor{white}{\textbf{Corrección}} \\
\midrule
\endhead
No aparecen filas con \lstinline|LIKE 'E%'| & No hay valores relacionados que
comiencen con E. & Quitar temporalmente el filtro y consultar
\lstinline|DISTINCT|. \\
\lstinline|invalid identifier| & Nombre de columna incorrecto, por ejemplo
\lstinline|DEPARTMENTS_ID|. & Revisar con \lstinline|DESCRIBE|; la columna es
\lstinline|DEPARTMENT_ID|. \\
\lstinline|column ambiguously defined| & Dos tablas poseen una columna con el
mismo nombre. & Anteponer el alias: \lstinline|e.department_id|. \\
Demasiadas filas o datos incoherentes & Falta una condición de unión o se unieron
columnas equivocadas. & Revisar cada \lstinline|JOIN ... ON| siguiendo claves
primarias y foráneas. \\
Una comparación con \lstinline|NULL| no funciona & Se escribió
\lstinline|= NULL|. & Usar \lstinline|IS NULL| o \lstinline|IS NOT NULL|. \\
Una fecha fija cambia de interpretación & Se confió en el formato regional. &
Usar \lstinline|DATE 'YYYY-MM-DD'| o \lstinline|TO_DATE| con formato explícito. \\
Error de agrupación & Una columna del \lstinline|SELECT| no está agrupada ni
dentro de una función de grupo. & Agregarla a \lstinline|GROUP BY| o aplicar una
función de grupo. \\
\bottomrule
\end{longtable}

\section{Método recomendado para construir una consulta}

Construir la consulta por capas permite localizar errores rápidamente y evita
escribir una sentencia extensa antes de comprobar sus relaciones básicas.

\begin{enumerate}
  \item Escribir con palabras qué datos debe mostrar el reporte.
  \item Identificar en qué tabla se encuentra cada columna.
  \item Elegir la tabla principal y escribir \lstinline|FROM|.
  \item Trazar el camino entre tablas y agregar cada \lstinline|JOIN ... ON|.
  \item Ejecutar sin filtros y revisar una muestra del resultado.
  \item Agregar \lstinline|WHERE| y comprobar los valores existentes con
        \lstinline|DISTINCT| cuando sea necesario.
  \item Agregar agrupaciones si el reporte usa totales o promedios.
  \item Ordenar al final con \lstinline|ORDER BY|.
  \item Reemplazar \lstinline|SELECT *| por las columnas exactas del reporte.
\end{enumerate}

\EndTheory
\begin{multicols}{3}[\section{Glosario alfabético de consulta rápida}]
\small
\raggedright
Esta sección concentra los comandos y conceptos de la unidad. Para aprender un
tema por primera vez, consulte la explicación y el ejemplo de su lección; use
este glosario para recordar su propósito durante la práctica.

\begin{glossarylist}
\item[ADD\_MONTHS] Suma meses a una fecha.
\item[AND] Exige que dos condiciones sean verdaderas.
\item[AS] Introduce un alias de columna; suele ser opcional.
\item[ASC / DESC] Define orden ascendente o descendente.
\item[AVG] Calcula un promedio.
\item[BETWEEN] Comprueba un rango inclusivo.
\item[CASE] Devuelve resultados diferentes según condiciones.
\item[COALESCE] Devuelve el primer valor no nulo.
\item[CONCAT] Une dos cadenas.
\item[COUNT] Cuenta filas o valores no nulos.
\item[CREATE TABLE] Crea una tabla.
\item[CROSS JOIN] Genera todas las combinaciones posibles entre dos tablas.
\item[DATE (literal ANSI)] Representa una fecha fija como
\lstinline|DATE 'YYYY-MM-DD'|, sin depender del formato de sesión.
\item[DECODE] Realiza comparaciones condicionales por igualdad en Oracle.
\item[DESCRIBE / DESC] Muestra la estructura de una tabla en el cliente Oracle.
\item[DISTINCT] Elimina duplicados del resultado.
\item[DUAL] Tabla especial de Oracle para evaluar expresiones.
\item[ESCAPE] Define el carácter que permite buscar \texttt{\%} o
\texttt{\_} literalmente en un patrón de \lstinline|LIKE|.
\item[FULL OUTER JOIN] Conserva filas coincidentes y no coincidentes de ambas tablas.
\item[FROM] Define el origen de datos.
\item[GROUP BY] Forma grupos para cálculos agregados.
\item[HAVING] Filtra grupos.
\item[IN] Compara con una lista de valores.
\item[INITCAP] Convierte la inicial de cada palabra a mayúscula.
\item[INNER JOIN] Devuelve únicamente las filas que coinciden en ambas tablas.
\item[INSERT INTO] Agrega una fila a una tabla.
\item[INSTR] Busca la posición de un texto dentro de otro.
\item[IS NULL] Comprueba si un valor es nulo.
\item[JOIN] Combina filas relacionadas de dos tablas.
\item[LAST\_DAY] Devuelve el último día del mes.
\item[LEFT JOIN] Conserva todas las filas de la tabla izquierda.
\item[LENGTH] Cuenta caracteres.
\item[LIKE] Compara un texto con un patrón.
\item[LPAD / RPAD] Completa un valor por la izquierda o la derecha.
\item[LOWER / UPPER] Convierte texto a minúsculas o mayúsculas.
\item[MAX / MIN] Obtiene el valor máximo o mínimo.
\item[MOD] Devuelve el resto de una división.
\item[MONTHS\_BETWEEN] Calcula meses entre dos fechas.
\item[NATURAL JOIN] Une columnas homónimas automáticamente; requiere especial cuidado.
\item[NEXT\_DAY] Devuelve el próximo día de semana indicado.
\item[NOT] Niega una condición.
\item[NOT BETWEEN / IN / LIKE] Niega el rango, la lista o el patrón correspondiente.
\item[NULL] Representa ausencia o desconocimiento de un valor.
\item[NULLIF] Devuelve nulo cuando dos expresiones son iguales.
\item[NVL / NVL2] Sustituye o clasifica valores nulos.
\item[ON] Declara la condición de un \lstinline|JOIN|.
\item[OR] Exige que al menos una condición sea verdadera.
\item[ORDER BY] Ordena el resultado final.
\item[REPLACE] Sustituye apariciones de un texto dentro de otro.
\item[RIGHT JOIN] Conserva todas las filas de la tabla derecha.
\item[ROUND / TRUNC] Redondea o trunca números y fechas.
\item[SELECT] Elige columnas y expresiones que se mostrarán.
\item[SUBSTR] Extrae una parte de una cadena.
\item[SUM] Suma los valores de un grupo.
\item[SYSDATE] Fecha y hora actuales del servidor Oracle.
\item[TO\_CHAR] Convierte números o fechas a texto con formato.
\item[TO\_DATE] Convierte texto a fecha con formato.
\item[TO\_NUMBER] Convierte texto a número.
\item[TRIM] Elimina espacios u otros caracteres en los extremos de un texto.
\item[USING] Abrevia una unión cuando la columna se llama igual en ambas tablas.
\item[WHERE] Filtra filas antes de agrupar.
\end{glossarylist}
\end{multicols}

\section*{Fuentes internas consultadas}
\addcontentsline{toc}{section}{Fuentes internas consultadas}

{\small
\textbf{Archivos SQL:}
\nolinkurl{U1/codigos/leccion-01.sql},
\nolinkurl{leccion-02.sql},
\nolinkurl{leccion-03.sql},
\nolinkurl{leccion-03-parte-02.sql},
\nolinkurl{leccion-04.sql},
\nolinkurl{leccion-04-parte-02.sql},
\nolinkurl{leccion-05.sql} y
\nolinkurl{leccion-06.sql}.
\textbf{Presentaciones:} lecciones 1 a 6 disponibles en
\nolinkurl{U1/presentaciones/}.\par}

\begin{ideaclave}[Antes de entregar una consulta]
Compruebe que las tablas y columnas existen, que cada \texttt{JOIN} sigue una
relación válida, que los textos y fechas usan literales correctos, que toda
columna no agregada figura en \texttt{GROUP BY} y que los alias describen con
claridad el resultado.
\end{ideaclave}

\end{document}
